\section{A positive magnetic mean and coercive energy}
\label{sec:magnetic-mean}

To estimate the contour sum, we center its ordinary particles by a
positive measure. A curved contour makes that measure real and positive.
The second part of the section proves that its real interaction controls
all species, uniformly when the contour approaches a boundary line.

\subsection{The contour carrying a positive magnetic mean}

From this point onward all real site offsets vanish.  Among the
ordinary \(X\) sites, a fraction \(t\in[1/2,1]\) has height
\(\beta\), and the rest have height zero.  All ordinary \(Y\) sites
have height zero, and either mark may have either height.  The contour
construction and the constants in the energy estimates do not depend on
the site counts; the centering identities hold exactly for each pair of
counts.  For the geometric comparison in
Figure~\ref{fig:magnetic-geometry}, write
\[
 L_X=\frac{\log s}{\gamma},\qquad L_Y=\frac{\log b}{\gamma}.
\]
These logarithmic wall intervals constrain the later particle estimates;
they are distinct from the contour carrying the mean.
\begin{figure}[H]
\centering
\begin{tikzpicture}[x=0.94cm,y=0.85cm,>=Latex,font=\small]
 \begin{scope}
  \draw[->] (-3.5,0)--(3.5,0) node[right] {$u$};
  \draw[densely dashed] (-3.2,1.05)--(3.2,1.05);
  \draw[densely dashed] (-3.2,2.1)--(3.2,2.1);
  \node[left] at (-3.2,1.05) {$\beta/2$};
  \node[left] at (-3.2,2.1) {$\beta$};
  \draw[blue!65!black,very thick,domain=-3.2:3.2,samples=80]
      plot (\x,{1.98-0.36*exp(-2.3*\x*\x)});
  \node[blue!65!black,anchor=west] at (0.7,1.35) {$z(u)=u+i\beta\theta(u)$};
  \node[below] at (0,0) {$0$};
  \draw[<->] (0,1.62)--(0,2.1);
  \node[above,anchor=south west] at (0.05,2.1) {$d_0$};
  \node at (0,-0.65) {Positive mean on the curved $X$ contour};
 \end{scope}
 \begin{scope}[yshift=-4.2cm]
  \fill[orange!20] (-3,0.15) rectangle (-1.8,0.62);
  \fill[blue!12] (-1.8,0.15) rectangle (1.8,0.62);
  \fill[orange!20] (1.8,0.15) rectangle (3,0.62);
  \draw[->] (-3.5,0)--(3.5,0) node[right] {$u$};
  \foreach \x/\lab in {-3/-L_Y,-1.8/-L_X,0/0,1.8/L_X,3/L_Y} {
    \draw (\x,-0.07)--(\x,0.09);
    \node[below] at (\x,-0.08) {$\lab$};
  }
  \node at (0,0.385) {common interval};
  \node[above] at (-2.4,0.67) {shoulder};
  \node[above] at (2.4,0.67) {shoulder};
  \draw[thick,blue!65!black] (-1.8,1.3)--(1.8,1.3);
  \draw[thick,orange!75!black] (-3,1.8)--(3,1.8);
  \node[right] at (1.85,1.3) {$X$ wall};
  \node[right] at (3.05,1.8) {$Y$ wall};
 \end{scope}
\end{tikzpicture}
\caption{Two distinct geometries in the magnetic estimate. Above, the
main $X$ contour carries a positive mean; the central distance $d_0$
can become small when almost all $X$ sites are upper sites. Below,
the particle numbers set different logarithmic wall lengths. The sharp
trial field treats the common interval and its two shoulders separately.
The curves and lengths are schematic.}
\label{fig:magnetic-geometry}
\end{figure}


The probability density and its transform are
\[
 \rho(w)=\frac{\sinh(2\gamma w)}
 {\sqrt3\,\beta\sinh(3\gamma w)},\qquad
 \widehat\rho(k)=\frac1{c_*(k)},\qquad
 c_*(k)=2\cosh(\beta k)-1 .
\]
The apparent singularity at \(w=0\) is removable.
Let \(\mathfrak R\) be its analytic primitive on
\(|\Im w|<\beta\), normalized by \(\mathfrak R(0)=0\).
We construct the contour as a level set of the imaginary primitive.
Its uniform Lipschitz bound will suffice for real-energy coercivity.
When \(t\) tends to one, the finer central estimates below will also
control smoothing and the particle coordinates used for fluctuations.

\begin{lemma}[Positive-density contour and its uniform regularity]
\label{lem:06-magnetic-curve}
For \(1/2\le t<1\), there is a unique real-analytic even function
\(\theta:\R\to[1/2,1)\) satisfying, with
\(z(u)=u+i\beta\theta(u)\),
\begin{equation}
 \Im\bigl[(1-t)\mathfrak R(z(u))
             +t\mathfrak R(z(u)-i\beta)\bigr]=0.
 \label{eq:06-curve}
\end{equation}
For \(t=1\), take \(\theta=1\) and \(\nu_X=\rho\).  For \(t<1\), define
\begin{equation}
 \nu_X(u)=
 [(1-t)\rho(z(u))+t\rho(z(u)-i\beta)]\,z'(u).
 \label{eq:06-density}
\end{equation}
In every case \(\nu_X\) is real, positive, even, and has integral one.
The graphs \(z\) are uniformly Lipschitz for \(t\in[1/2,1]\).
There is a unique \(\theta_\infty\in[1/2,1]\) satisfying
\begin{equation}
 (1-t)\sin(\gamma\beta\theta_\infty)
 +t\sin(\gamma\beta(\theta_\infty-1))=0,
 \label{eq:06-limit-height}
\end{equation}
and \(\theta(u)-\theta_\infty\) tends to zero exponentially with every
fixed-order derivative, uniformly in \(t\).
Outside any fixed neighborhood of zero, the curves and the
correspondingly rescaled tail densities
\(e^{\gamma|u|}\nu_X(u)\) have uniform analytic bounds, and
\(\nu_X(u)\asymp e^{-\gamma|u|}\).

More precisely, put \(r=1-t>0\),
\(\Delta(u)=\beta(1-\theta(u))\), and \(d_0=\Delta(0)\).
As \(r\downarrow0\),
\begin{equation}
 d_0\asymp r\log(1/r),\qquad
 \sup_u\left(
 |z(u)-(u+i\beta)|+|z'(u)-1|
 +\left|\frac{\nu_X(u)}{\rho(u)}-1\right|\right)=o(1).
 \label{eq:06-center-regularity}
\end{equation}
On a fixed small interval around zero, for each integer \(j\ge1\),
\[
 (|u|+d_0)^j
 \left(
 |\partial_u^j(z'(u)-1)|
 +|\partial_u^j(\nu_X(u)-\rho(u))|
 \right)=o(1)
\]
uniformly in \(u\).
The odd increasing change of coordinates \(b_X\) defined by
\[
 \int_0^{b_X(u)}\rho(v)\,dv=\int_0^u\nu_X(v)\,dv
\]
has \(b_X(u)-u=o(1)\) and \(b_X'(u)=1+o(1)\) uniformly as
\(t\uparrow1\), with the same scaled estimates for positive-order
derivatives of \(b_X'-1\).
\end{lemma}

\begin{proof}
The choice at \(t=1\) has all the asserted positivity and mass properties.
We construct the contour for \(t<1\), then prove the uniform estimates
as \(t\uparrow1\).
First \(\Re\rho(w)>0\) throughout \(|\Im w|<\beta\).
To verify the sign away from the removable zero, write
\(\gamma w=x+iy\), so \(|y|<\pi/3\).
The numerator of the real part, up to a positive factor, is
\[
 \cosh(5x)\cos y-\cosh x\cos(5y).
\]
Here \(\cos y>0\) and
\(\cos y-\cos(5y)=2\sin(3y)\sin(2y)\ge0\);
strict positivity follows except at the removable case \(x=y=0\),
where \(\rho(0)>0\).

Set
\(\Phi_t(u,y)=\Im[(1-t)\mathfrak R(u+iy)
 +t\mathfrak R(u+i(y-\beta))]\).
Its derivative in \(y\) is strictly positive for \(0<y<\beta\).
By conjugation symmetry,
\(\Phi_t(u,\beta/2)\le0\) for \(t\ge1/2\).
At \(y\uparrow\beta\), its second summand tends to zero and the
first is strictly positive, possibly infinite at \(u=0\).
The intermediate value theorem and strict monotonicity give the
unique root.  The analytic implicit-function theorem gives
real analyticity, and symmetry in \(u\) gives evenness.

For positivity of \(\nu_X\), write
\(A=(1-t)\rho(z)+t\rho(z-i\beta)=a+ib\), with \(a>0\).
Differentiating \(\Phi_t(u,\beta\theta(u))=0\) gives
\(\beta\theta'=-b/a\), and hence
\[
 \nu_X=A(1+i\beta\theta')=\frac{|A|^2}{a}>0.
\]
Evenness follows from conjugation symmetry.
Moving each translated density contour within its pole-free strip
shows \(\int\nu_X=1\).  Equivalently the primitive in
\eqref{eq:06-curve} tends to \(1/2\) and \(-1/2\) at its two ends.

The exponential expansion of \(\rho\) at either end begins with a
positive multiple of \(e^{-\gamma|u|}\).
Substitution into the imaginary primitive equation, divided by
this leading term, gives \eqref{eq:06-limit-height}.
Its derivative with respect to \(\theta_\infty\) is bounded below
uniformly, since both relevant cosine arguments lie in
\([-\pi/3,\pi/3]\).  The analytic implicit-function theorem applied
to the exponentially small remainder proves the uniform tail
statements, including fixed-order derivatives.  On compact
\(t\)-subintervals below one, the same conclusions in the remaining
bounded \(u\)-region follow by compactness and the strict derivative
in \(y\).

It remains to examine the possible central degeneration as
\(r=1-t\downarrow0\).  At \(w=i\beta\), the density has a simple
pole with negative imaginary residue, explicitly
\(-i/(6\beta\gamma)\).
On a sufficiently small fixed interval in \(u\), the defining
equation therefore compares
\[
 -(1-r)\Im\mathfrak R(u-i\Delta)\asymp\Delta
\]
with
\[
 r\Im\mathfrak R(u+i\beta-i\Delta)
 \asymp r\left(1+\log^+\frac1{|u|+\Delta}\right).
\]
At \(u=0\) this gives \(d_0\asymp r\log(1/r)\).
The same estimates, separating \(|u|\le d_0\) and \(|u|>d_0\),
give \(|u|+\Delta(u)\asymp |u|+d_0\).

To obtain uniform derivative bounds, write the defining equation as
\[
 F_r(u,\Delta)
 =r\Im\mathfrak R(u+i\beta-i\Delta)
  +(1-r)\Im\mathfrak R(u-i\Delta)=0.
\]
The regular function vanishes identically when \(\Delta=0\), and its
Taylor expansion on this fixed central interval is
\[
 \Im\mathfrak R(u-i\Delta)=-\Delta\rho(u)+O(\Delta^3).
\]
Thus its \(u\)-derivative is \(O(\Delta)\), whereas its
\(\Delta\)-derivative is \(-\rho(u)+O(\Delta^2)\).
Either first derivative of the pole summand, including its coefficient,
is bounded by
\[
 O\!\left(\frac{r}{|u|+d_0}\right)
 =O\!\left(\frac1{\log(1/r)}\right)=o(1).
\]
Since \(\rho\) is bounded below here and \(\Delta=o(1)\), it follows that
\[
 |\partial_\Delta F_r|\ge c_3>0,\qquad
 |\partial_u F_r|
 \le C\left(\Delta+\frac{r}{|u|+d_0}\right)=o(1).
\]
In particular \(\Delta'=-\partial_u F_r/\partial_\Delta F_r=o(1)\).
For higher derivatives, extend each imaginary part analytically using
the two conjugate analytic terms.  Around a real solution at \(u_0\),
take complex discs in \(u,\Delta\) of radii
\(c_2(|u_0|+d_0)\).  The pole remains outside these discs when
\(c_2\) is sufficiently small.  The \(\Delta\)-derivative remains
bounded away from zero, while varying \(u\) with \(\Delta=\Delta(u_0)\)
changes \(F_r\) by \(o(1)(|u_0|+d_0)\), by the same estimates.
The analytic implicit root therefore extends to a disc of comparable
radius with oscillation \(o(1)(|u_0|+d_0)\).  Cauchy's estimates give
\(z'-1=o(1)\) and the stated scaled higher-derivative bounds.
The formula for \(\nu_X\) gives the same estimates for
\(\nu_X-\rho\).  In this central region \(\rho\) is bounded below;
off it the uniform tail expansion supplies the relative estimate
in \eqref{eq:06-center-regularity}.  These statements prove, in
particular, uniform Lipschitzness and uniform comparison
\(\nu_X\asymp\rho\) throughout the parameter interval.

Finally the cumulative distributions are strictly increasing and
have identical total mass.  Analytic inversion on bounded
intervals, and inversion of their normalized exponential tails,
show that their matching map \(b_X\) is odd and obeys the claimed
uniform displacement and derivative bounds.  Differentiating
\(\rho(b_X)b_X'=\nu_X\) on the same analytic discs gives the scaled
higher-derivative estimates.
\end{proof}

\subsection{Exact centering of the ordinary and marked factors}

Label the six species by \(I=(\sigma,J)\), where \(J=0,F,P\)
denotes main, far singleton, or far double, and let
\(w_I=1,1,2\), respectively.  On each species use the real coordinate
\(u\) as its integration coordinate; only the main \(X\) species
follows the curved graph \(z(u)\).
Let \(H\) be minus the sign-normalized particle--particle logarithm,
and \(V\) the sign-normalized ordinary site--particle logarithm,
with the following logarithm convention. Reduce each angle to
$[0,2\pi)$ and use linear combinations of
\[
 K_\vartheta(u)=-\log\!\left[2\left(\sin(\vartheta/2)\cosh(u/2)
       -i\cos(\vartheta/2)\sinh(u/2)\right)\right],\quad
 K_0(u)=-\log|2\sinh(u/2)|.
\]
The particle factors determine the exponents in $H$; the negative
site--particle exponents determine $V$. Reverse cross entries by
transposing and reflecting the argument. Fourier transforms use
$\widehat f(k)=\int e^{-iku}f(u)\,du$ and inverse measure $dk/(2\pi)$.
At nonzero frequency,
\[
 \widehat K_\vartheta(k)=\frac{\pi e^{(\pi-\vartheta)k}}
                              {k\sinh(\pi k)},\qquad
 \widehat K_0(k)=\frac{\pi\cosh(\pi k)}{k\sinh(\pi k)}.
\]
The transforms follow by differentiating twice and taking the double
pole residue; symmetric averaging gives the angle-zero case.
Affine ambiguities are fixed by the real tails $-|u|/2+O(e^{-|u|})$
and the odd imaginary tails
$(\pi-\vartheta)\operatorname{sign}(u)/2$. The forms extend
bilinearly to complex measures.

For the curved main self interaction, retain the horizontal
logarithm and add the continuous logarithm of the ratio of its
curved and horizontal factors.  In particular the two vanishing
sinh factors are divided before this correction is taken.
Choose the logarithm that is zero before contour deformation.
This prescription fixes the diagonal continuation and gives kernels
conjugated by simultaneous negation of both real coordinates.
Constant signs depend only on the species counts.

Put
\[
 \mu_X=s\nu_X(u)\,du,\qquad
 \mu_Y=b\rho(u)\,du,\qquad \mu_I=0\quad(I\ {\rm far}),
\]
and write \(\mathfrak t_I^{\,a}\) for the particle logarithm of the
single-mark factor
\(\prod_vP_{-4\sigma_v}(v-a)/P_{\sigma_v}(v-a)\).

\begin{lemma}[Magnetic mean identities]
\label{lem:06-magnetic-mean}
The ordinary one-body logarithm is \(H\mu\).
Multiplying the ordinary site-pair factors by
\(\exp((\mu,H\mu)/2)\) gives, up to a constant sign,
\begin{equation}
 C_*^n\prod_{\substack{i<j\\i,j\in X\cup Y}}H_*(a_i-a_j)^2,
 \qquad
 C_*=\frac{(1+d)\gamma}{4d\cos(\beta/2)}.
 \label{eq:06-ordinary-mean}
\end{equation}
For each marked log \(a\), its site factors in \(\mathcal T\),
multiplied by \(\exp((\mathfrak t^{\,a},\mu))\), give, again up to
a constant sign,
\begin{equation}
 \prod_{i\in X\cup Y}H_*(a_i-a)^2.
 \label{eq:06-marked-mean}
\end{equation}
On either main contour the real part of
\(\mathfrak t_I^{\,a}\), together with its first derivative, decays
exponentially at infinity.  Its imaginary part tends exponentially
to
\(\sigma_I(3\beta/2)\operatorname{sign}u\).
The logarithm is regular with uniformly bounded first derivative
at the center.  On far species one has an exponentially decreasing
upper bound for its real part; zeros of its numerator are allowed.
\end{lemma}

\begin{proof}
We justify the contour continuations in the Fourier mean
identities, since the \(X\) site measure now has two different
heights.  For a test on the \(X\) main contour, slide separately the
two translated \(\rho\) contours to a horizontal line through the
test point, keeping that line at distance strictly less than
\(\beta\) from the relevant site height.  For the diagonal
logarithm, split the contour at the test point and deform its two
halves with this common endpoint.  Its logarithmic singularity is
integrable, and the chosen continuous branch crosses no cut.
For any other test species use a fixed admissible horizontal
height between \(\beta/2\) and \(\beta\).

No additional singularity is met in these slides.
The signed same-group separation of a far elementary particle from
a main contour or its associated horizontal site line has absolute
value strictly between \(2\beta\) and \(5\beta\).
The cross separation, with the \(X\) coordinate first, is positive
and less than \(6\beta\).
All these bounds follow directly from the listed contours.
They also hold for the slides used in a second pairing with a
site's translated density.  A same-group \(V\) is always slid
within distance \(\beta\) of its site.  Coincident site--main
heights in a cross \(V\) are regular.

On horizontal contours Fourier transforms use
\(\widehat f(k)=\int e^{-iku}f(u)\,du\); put \(x=\beta k\).
The self-main entry has transform
\[
 \widehat H_{\rm self}(k)
 =\frac{2\pi c_*(k)\cosh(3x)}{k\sinh(\pi k)}.
\]
For a same-group far elementary particle at signed height
separation \(j\beta\), replace the numerator
\(c_*\cosh(3x)\) by
\[
 c_*\cosh(x)e^{(4\operatorname{sign}j-j)x}.
\]
For a cross elementary pair at separation \(d\beta>0\), replace it
by \(c_*e^{(3-d)x}\), transposing by \(k\mapsto-k\) when the row
and column are reversed.
Convolution with each translated \(\rho\) cancels \(c_*\);
the translation multiplier adjusts the height from its main
contour to its site.  These are exactly the transforms of the
corresponding \(V\) entries.  This proves the one-body identity.

Pair once more with the translated density of a second site.
For a horizontal same-group pair the symbol of
\(\log B_s+\rho*H_{\rm self}*\rho\), divided by
\(2\pi/(k\sinh(\pi k))\), is
\[
 -\cosh(2x)-\cosh x+\frac{\cosh(3x)}{c_*(k)}
 =-2\cosh x+\cosh(4x)-\sinh(4x)\tanh(3x/2).
\]
The expression on the right is the symbol of \(\log H_*^2\) in the
same normalization: factor \(H_*\) and use
\(\widehat{\log|\tanh(\gamma u/2)|}(k)
=-\pi\tanh(3\beta k/2)/k\).
This is the horizontal centering calculation of
\cite[Lemma~3.1]{Electric}; the contour slides above extend it to
the two site heights here.  In the cross case one may first compute
at positive main--main separation and undo that shift in the
transforms.  The cancellation needed there is
\[
 \frac{\sinh(3x)}{c_*(k)}=\sinh(2x)+\sinh x.
\]
Endpoint identities at a height difference \(\beta\) follow by
moving the heights slightly inward and taking the regular limit.
The half self term for each site is fixed by the same-group identity
at zero.  With the affine terms fixed by the tails as below, this gives
\[
 \exp\!\left(\tfrac12(\rho*H_{\rm self}*\rho)(0)\right)
 =\frac{H_*(0)}{\sqrt{B_s(0)}}
 =\frac{\gamma(1+d)}{4d\cos(\beta/2)}=C_*.
\]
Here \(H_*(0)=\gamma(1+d)\) and
\(\sqrt{B_s(0)}=4\sin\beta\sin(3\beta/2)=4d\cos(\beta/2)\),
with the positive square root.  There is one such factor per ordinary
site, proving \eqref{eq:06-ordinary-mean} with its exact normalization.

For a single mark, slide the density in the same way.
After removing the ordinary-site to marked-site translation
multiplier, the combined log symbol, divided by
\(-\pi/(k\sinh(\pi k))\), is
\[
 e^{2\sigma x}+e^{\sigma x}
 +\frac{1-e^{3\sigma x}}{c_*(k)}
 =
 2\left(\cosh(2x)+\cosh x-\frac{\cosh(3x)}{c_*(k)}\right).
\]
The equality follows by multiplying by \(2\cosh x-1\).
This is exactly the symbol of \(\log H_*^2\), proving
\eqref{eq:06-marked-mean}.
These transform comparisons determine the second derivatives.
To fix the affine ambiguity, subtract the linear real tails:
the remaining real constants vanish by the exponential expansions
and the centered real offsets; a translated \(\rho\) has purely
imaginary first moment.  On horizontal contours the imaginary
parts are bounded and odd, so their affine ambiguity vanishes as
well.  Continuing the chosen logarithms fixes the stated signs.

Finally on a main contour the numerator angle of a marked particle
factor is reduced to \(4\) plus the relevant shift, in units of
\(\beta\), and the denominator angle is \(1\) or \(7\) plus that
shift.  These angles stay in the open interval \((0,8)\).
The log-sinh tail expansion therefore gives zero real limit,
exponential real decay with a derivative, and the imaginary
limits \(\sigma(3\beta/2)\operatorname{sign}u\).
The compact part remains uniformly regular because the main
angles stay away from zeros and \(z'\) is uniformly bounded by
Lemma \ref{lem:06-magnetic-curve}.
On far contours the denominator remains regular and the same
tail expansion bounds the log modulus above by \(Ce^{-c_2|u|}\).
A numerator zero only improves this upper bound.
\end{proof}

\subsection{Positivity of the real energy}

For a real vector of smooth rapidly decreasing densities
\(\xi=(\xi_I)\), assume weighted neutrality and define
\[
 \int_{\R}\sum_Iw_I\xi_I=0,\qquad
 d_\xi=\sum_Iw_I\xi_I,\qquad
 D(u)=\int_{-\infty}^u d_\xi(v)\,dv,\qquad
 B(\xi)=\tfrac12(\xi,\Re H\,\xi).
\]
Use the inhomogeneous negative Sobolev norm
\[
 \|\xi_I\|_{-1/2}^2
 =\int_{\R}(1+k^2)^{-1/2}|\widehat\xi_I(k)|^2
 \,\frac{dk}{2\pi}.
\]

\begin{theorem}[Uniform real-energy coercivity]
\label{thm:06-real-positivity}
For the magnetic contours there are constants \(0<c_2<C_2<\infty\)
such that
\begin{equation}
 c_2\left(\|D\|_2^2+\sum_I\|\xi_I\|_{-1/2}^2\right)
 \le B(\xi)\le
 C_2\left(\|D\|_2^2+\sum_I\|\xi_I\|_{-1/2}^2\right).
 \label{eq:06-real-positivity}
\end{equation}
They are uniform in \(t\in[1/2,1]\) and in all sufficiently small
positive far-contour parameters \(\varepsilon\).
The form extends to the completion in this norm.
\end{theorem}

\begin{proof}
We will sweep the curved \(X\) charge harmonically to the two edges
of the strip between heights \(\beta/2\) and \(\beta\).
The boundary energy and the remaining positive Green energy together
will control both the species norms and the weighted total primitive.

\par\smallskip\noindent\textbf{The horizontal boundary lines.}
Allow independent \(X\) main charges on the two boundary lines.
This gives a translation-invariant system with seven species.
At \(\varepsilon=0\), their elementary
height lists, in units of \(\beta/8\), are
\[
 (4),\quad(8),\quad(25),\quad(24,32),\quad
 (0),\quad(-28),\quad(-24,-32).
\]
The first four species belong to \(X\); the final three belong to
\(Y\).  Divide the real Fourier symbol by
\(\pi/(k\sinh(\pi k))\), and divide its \(IJ\) entry by
\(w_Iw_J\).  Expand the resulting matrix in
\(v=\cosh(\beta k/8)-1\).  Its constant matrix has every entry
equal to \(2\), and every coefficient of degree
\(1\le l\le32\) is strictly positive definite.
Here is a finite rational verification of this assertion.

The coefficient of \(v^l\) in \(\cosh(32\,\operatorname{arcosh}(1+v))\)
is
\[
 c_{32,l}=\frac{2^l}{(2l)!}
 \prod_{a=0}^{l-1}(32^2-a^2)>0.
\]
After division by this coefficient, define
\[
 h_l(j)=
 \prod_{a=0}^{l-1}\frac{i_j^2-a^2}{32^2-a^2},
 \qquad i_j=(j\bmod64)-32.
\]
For each pair of elementary heights, let \(j\) be their absolute
difference.  Average over the two species' height lists the
expression
\[
 2h_l(j)+h_l(j+16)+h_l(j-16)-h_l(j-8)-h_l(j+8)
\]
when they are in the same group, and
\[
 2\bigl(h_l(j)+h_l(j+16)-h_l(j+8)\bigr)
\]
when they are in opposite groups.  Denote the resulting symmetric
matrix by \(M_l\).
These formulas follow by summing the log-factor angles modulo
\(2\pi=8\beta\).  The elementary cosh coefficient formula follows
from
\((2v+v^2)y''+(1+v)y'=i_j^2y\).

For \(l\le12\), symmetric elimination gives successive pivots
strictly larger than the rational decimals in the following table:
\[
\begin{array}{c|rrrrrrr}
 l&1&2&3&4&5&6&7\\ \hline
 1&1&.7&.8&.27&.49&.59&.09\\
 2,3&1&1&1&.32&.8&1&.30\\
 4,5,6&1.7&1.65&1.48&.35&1.53&1.65&.45\\
 7,\ldots,12&1.9&1.88&1.74&.35&1.89&1.88&.48
\end{array}
\]
The entries are obtained using exact rational substitution in
\[
 m_{ij}^{(p)}
 =m_{ij}^{(p-1)}
 -\frac{m_{ip}^{(p-1)}m_{jp}^{(p-1)}}{m_{pp}^{(p-1)}},
 \qquad i,j>p,\qquad M_l=(m_{ij}^{(0)}).
\]
Thus all these matrices are positive definite.

For \(13\le l\le32\), each \(h_l(j)\) is nonnegative and
nonincreasing in \(l\): before a zero occurs, every additional
factor lies in \([0,1]\), and after that zero the product stays
zero.  Split the displayed averages into their positive and
negative parts, \(M_l=P_l-Q_l\).  A lower bound for its diagonal
is \(\operatorname{diag}(P_{32}-Q_{13})\), and an entrywise upper
bound for the absolute off-diagonal entries is
\[
 \max(P_{13}-Q_{32},\,Q_{13}-P_{32}).
\]
Weight column \(j\) relative to row \(i\) by \(v_j/v_i\), with
\(v=(1,1,1,2,1,1,1)\).
Subtracting the weighted off-diagonal row sums from the diagonal
lower bounds gives exactly
\[
\begin{split}
 (&539085539017/288991072832,\ 
 153103458093/288991072832,\ 
 6754920483/37289170688,\\
 &2413205499/288991072832,\ 
 8798842243/4515485513,\\
 &496150561213/288991072832,\ 
 32952783649/1155964291328).
\end{split}
\]
All seven numbers are positive.  Weighted strict diagonal
dominance of the symmetric \(M_l\) proves positive definiteness
for every remaining degree.

At small frequencies the constant coefficient supplies the
rank-one \(k^{-2}\) term for the weighted total density, and the
strictly positive first coefficient supplies a uniformly positive
bounded term on all species.  At large frequencies the top
coefficient supplies a positive multiple of \(1/|k|\), since
\(\pi=4\beta\).  The intervening frequencies are controlled by
continuity and positivity.  The real tail
\(-w_Iw_J|u-v|\) contributes exactly \(\|D\|_2^2\) to \(B\).
These observations prove the two-sided norm estimate on the
seven boundary-line species at \(\varepsilon=0\).

\par\smallskip\noindent\textbf{Uniform movement of the far double.}
The estimate continues uniformly to sufficiently small
\(\varepsilon>0\).  At small frequencies use the strict positive
quadratic coefficient, and on compact frequency intervals use
continuity.  To check uniformity at infinite frequency, divide
the normalized symbol further by \(\cosh(4\beta k)\).
For \(0\le\varepsilon<1/16\) it has the form
\[
 \operatorname{diag}(2,2,2,\tfrac12,2,2,\tfrac12)
 +q_\varepsilon(k)(E_{24}+E_{42})
 +O(e^{-c_3|k|}),\qquad
 q_\varepsilon(k)=
 \frac{\cosh((4-\varepsilon)\beta k)}{2\cosh(4\beta k)}.
\]
Here \(E_{ij}\) is the elementary matrix, and the error is uniform
in \(\varepsilon\): after collecting the elementary factors of
each double, all other contributing angles stay a fixed distance
from zero modulo \(8\beta\).  In particular the potentially small
angles in the \(X\)-double/\(Y\)-double entry cancel with total
coefficient \(2-4+2=0\), for every \(\varepsilon\).
The exceptional pair couples the \(X\) main line of height
\(\beta\) to the \(X\) double, and \(0\le q_\varepsilon\le1/2\).
The only nontrivial \(2\) by \(2\) block is
\(\left(\begin{smallmatrix}2&q\\q&1/2\end{smallmatrix}\right)\);
its least eigenvalue is at least
\((5-\sqrt{13})/4>0\).
This proves uniform coercivity also in every joint limit
\(\varepsilon\downarrow0\), \(|k|\to\infty\).

\par\smallskip\noindent\textbf{Harmonic sweeping to the boundary.}
We now pass from the two boundary lines to the graph.
Let \(\mathcal S=\{z:\beta/2<\Im z<\beta\}\), and balayage the
real charge on the graph harmonically to the two edges of
\(\mathcal S\).  Leave all other species unchanged, and denote the
resulting seven-species charge by \(\widetilde\xi\).
All interactions are harmonic in this strip, in either variable
away from the source, and have at most linear growth.
The only interior source in a graph self interaction is
\[
 -2\log\left|2\sinh\frac{z-z'}2\right|,
\]
whose logarithmic singularity is \(-2\log|z-z'|\).
Let \(G_{\mathcal S}\) be the Dirichlet Green kernel of
\(-\Delta\), with singularity
\(-(2\pi)^{-1}\log|z-z'|\), and put
\[
 \mathcal G(\xi_X)=
 \iint G_{\mathcal S}(z(u),z(v))\,\xi_X(u)\xi_X(v)\,du\,dv.
\]
Harmonic interpolation first in one variable and then in the other
therefore gives the exact identity
\begin{equation}
 B(\xi)=B_{\partial\mathcal S}(\widetilde\xi)
              +2\pi\mathcal G(\xi_X).
 \label{eq:06-balayage}
\end{equation}
The coefficient \(2\pi\) includes the factor \(1/2\) in the
definition of \(B\).
The Green energy is nonnegative for signed real charges: by
Dirichlet inversion and integration by parts it is the squared
dual norm of the charge acting on traces of zero-boundary
functions with gradient \(L^2\)-norm at most one.
Harmonic measures in the strip decay exponentially; their Fourier
kernels are ratios of the form
\(\sinh(yk)/\sinh(\beta k/2)\).
These facts justify the interpolations for rapidly decreasing
densities and then by completion.

\par\smallskip\noindent\textbf{Uniform control as the graph approaches an edge.}
For clarity, we verify that \eqref{eq:06-balayage} gives exactly
the claimed norm, even when the graph approaches the upper edge.
The Sobolev trace principle underlies this comparison
\cite[Theorem~1.I]{Gagliardo1957Trace}; we give the strip construction
to retain uniform constants as the graph approaches an edge.
The trace map \(H^1(\R^2)\to H^{1/2}(\R)\) on the graph and a right
inverse both have bounds depending only on its uniform Lipschitz
constant.  Flatten the graph by a uniformly bilipschitz map.
On a horizontal line, the trace bound follows by transverse
Fourier Cauchy--Schwarz, and a right inverse extends a mode \(k\)
by \(e^{-\sqrt{1+k^2}|y|}\).
Reflection across the strip edges with a cutoff gives the same
trace bound for strip functions.
Harmonic extension from the two boundaries is bounded
\(H^{1/2}\to H^1\); this follows either from the displayed
Fourier kernels or the Dirichlet principle together with the
fixed strip width.

Duality now bounds the two projected charges' \(H^{-1/2}\) norms
by the graph charge's norm.  It also bounds the Green energy above
by that squared norm, since zero-boundary strip functions have
their full \(H^1\) norm controlled by their gradient.
Conversely, extend an arbitrary \(H^{1/2}\) test on the graph to
the plane, restrict it to the strip, and decompose it into its
harmonic boundary extension and a zero-boundary remainder.
Pairing these two terms with the charge gives
\[
 \|\xi_X\|_{-1/2}^2
 \le C_3\left(
 \|\widetilde\xi_{X,\beta/2}\|_{-1/2}^2+
 \|\widetilde\xi_{X,\beta}\|_{-1/2}^2+
 \mathcal G(\xi_X)\right).
\]
Thus the Sobolev parts of the norms are equivalent in both
directions, with the Green term included.

It remains to compare primitives of the weighted total densities.
Let \(\widetilde D\) be the primitive after projection.
For a smooth compactly supported function \(b\), pair the
horizontal marginal of the difference between graph charge and
projected charge with \(b\).  This pairing equals that of the
graph charge with \(b(u)\) minus the harmonic extension of the
same boundary value \(b\) on both edges.
The latter difference has zero boundary values and gradient norm
at most \(C_4\|b'\|_2\): compare the harmonic extension to the
constant-in-height function \(b(u)\) by the Dirichlet principle.
Green-energy duality, followed by integration by parts in \(u\),
therefore yields
\[
 \|D-\widetilde D\|_2^2\le C_5\mathcal G(\xi_X).
\]
Together with \eqref{eq:06-balayage}, the boundary-line coercivity,
and the Sobolev norm comparisons, this proves both inequalities
in \eqref{eq:06-real-positivity}.
All constants depend only on the fixed strip width and uniform
Lipschitz bounds, so remain uniform as the graph reaches an edge,
including the case \(t=1\).
\end{proof}
